\section{Conclusion}
\label{sec:conclusion}

We have mechanized the connection between a source semantics and a
Goblint-style, context-sensitive, side-effecting equation system solved by a
verified solver, with activation traces as the concrete counterpart of
contexts. The theorem is a first step. The result is partial correctness, VIMP is
scalar, the theorem is per program point, and precision is shown only by
witnesses for single programs.

\subparagraph*{Per-context certification.}
\Cref{thm:routed} bounds each context by its solved state, but the report-level
theorem keeps only the union over contexts. Proving that the report's row for
a context covers $\ActColl(v, c)$ would certify the per-context states.

\subparagraph*{Memory.}
The contract refers to the language only through a few step functions, so
features such as machine integers leave it intact. The trace semantics,
however, assumes that a callee cannot modify its caller's locals. Pointers into
stack frames and threads break this assumption; a heap without such pointers
needs only a richer store model.

\subparagraph*{Termination.}
Known termination results for top-down solvers assume finitely many
unknowns~\cite{tilscher26jar}. Entry-state contexts need not be finite, and
all timeouts in \cref{sec:evaluation} come from them. Goblint bounds them with
optional widening on contexts (\texttt{ana.context.widen}) or a context gas; a
general result would need such a bound in the routing.
